\documentclass[11pt,a4paper]{article}
\usepackage{amsmath,amssymb}
\usepackage{listings}
\usepackage{hyperref}
\usepackage[margin=1in]{geometry}
\usepackage{fancyhdr}

\pagestyle{fancy}
\fancyhf{}
\fancyhead[L]{Module 04: Functions}
\fancyhead[R]{C Programming Course}
\fancyfoot[C]{\thepage}

\lstset{
  language=C,
  basicstyle=\ttfamily\small,
  keywordstyle=\bfseries,
  commentstyle=\itshape,
  numbers=left,
  numberstyle=\tiny,
  frame=single,
  breaklines=true,
  tabsize=4
}

\title{Module 04: Functions}
\author{C Programming Course}
\date{}

\begin{document}
\maketitle
\thispagestyle{fancy}

\section{Function Declaration and Definition}
A function declaration (prototype) tells the compiler about the function signature.
The definition provides the body.

\begin{lstlisting}
#include <stdio.h>

/* Declaration */
int add(int a, int b);

int main(void) {
    printf("3 + 4 = %d\n", add(3, 4));
    return 0;
}

/* Definition */
int add(int a, int b) {
    return a + b;
}
\end{lstlisting}

\section{Function Parameters and Return Types}
C passes arguments by value. To modify the caller's variable, pass a pointer.

\begin{lstlisting}
void swap(int *a, int *b) {
    int temp = *a;
    *a = *b;
    *b = temp;
}
\end{lstlisting}

\section{Scope and Lifetime}
\begin{itemize}
  \item \textbf{Local variables} -- visible only within their enclosing block.
  \item \textbf{Global variables} -- visible across the entire file.
  \item \textbf{Static local variables} -- retain their value between calls.
\end{itemize}

\begin{lstlisting}
#include <stdio.h>

void counter(void) {
    static int count = 0;
    count++;
    printf("Called %d times\n", count);
}

int main(void) {
    counter();
    counter();
    counter();
    return 0;
}
\end{lstlisting}

\section{Recursion}
A function that calls itself. Every recursive function needs a base case.

\begin{lstlisting}
int factorial(int n) {
    if (n <= 1) return 1;
    return n * factorial(n - 1);
}
\end{lstlisting}

\section{Function Pointers}
A pointer that stores the address of a function.

\begin{lstlisting}
#include <stdio.h>

int square(int x) { return x * x; }

int main(void) {
    int (*fp)(int) = square;
    printf("5 squared = %d\n", fp(5));
    return 0;
}
\end{lstlisting}

\end{document}
